\documentclass[aspectratio=169]{beamer}
\usetheme{Madrid}
\usecolortheme{default}
\usepackage[utf8]{inputenc}
\usepackage[spanish]{babel}
\usepackage{graphicx}
\usepackage{hyperref}

\title{Normas APA - Séptima Edición}
\subtitle{Lista de Referencias}
\author{Edison Achalma}
\institute{Universidad Nacional de San Cristóbal de Huamanga}
\date{\today}

\begin{document}

\frame{\titlepage}

\begin{frame}{Tabla de Contenido}
\tableofcontents
\end{frame}

\section{¿Qué son las Referencias?}

\begin{frame}{Definición}
\begin{block}{Referencias}
Listado con la información completa de todas las fuentes citadas en el texto
\end{block}

\vspace{0.5cm}
\textbf{Funciones:}
\begin{itemize}
    \item Dar crédito a los autores
    \item Permitir localización de las fuentes
    \item Facilitar verificación de la información
    \item Demostrar rigor académico
\end{itemize}
\end{frame}

\begin{frame}{Referencias vs Bibliografía}
% Imagen: Tabla comparativa entre referencias y bibliografía
\begin{center}
\includegraphics[width=0.85\textwidth]{images/referencias_vs_bibliografia.png}
\end{center}

\begin{alertblock}{En APA}
Usamos REFERENCIAS, no bibliografía
\end{alertblock}
\end{frame}

\begin{frame}{Formato General}
% Imagen: Ejemplo de página de referencias con elementos señalados
\begin{center}
\includegraphics[width=0.75\textwidth]{images/formato_general_referencias.png}
\end{center}
\end{frame}

\section{Elementos de una Referencia}

\begin{frame}{Los Cuatro Elementos Básicos}
% Imagen: Diagrama mostrando Autor - Fecha - Título - Fuente
\begin{center}
\includegraphics[width=0.85\textwidth]{images/cuatro_elementos_referencia.png}
\end{center}
\end{frame}

\begin{frame}{Elemento 1: Autor}
% Imagen: Ejemplos de formato de autor según número (1, 2, 3-20, 21+, sin autor)
\begin{center}
\includegraphics[width=0.9\textwidth]{images/elemento_autor.png}
\end{center}
\end{frame}

\begin{frame}{Elemento 2: Fecha}
% Imagen: Ejemplos de diferentes formatos de fecha (año, mes y año, sin fecha)
\begin{center}
\includegraphics[width=0.85\textwidth]{images/elemento_fecha.png}
\end{center}
\end{frame}

\begin{frame}{Elemento 3: Título}
% Imagen: Ejemplos de títulos de libro, artículo, capítulo con formato correcto
\begin{center}
\includegraphics[width=0.85\textwidth]{images/elemento_titulo.png}
\end{center}
\end{frame}

\begin{frame}{Elemento 4: Fuente}
% Imagen: Ejemplos de fuente para libro (editorial), artículo (revista), web (URL)
\begin{center}
\includegraphics[width=0.85\textwidth]{images/elemento_fuente.png}
\end{center}
\end{frame}

\section{Referencias de Libros}

\begin{frame}{Libro con Autor}
% Imagen: Ejemplo completo de referencia de libro con autor
\begin{center}
\includegraphics[width=0.85\textwidth]{images/referencia_libro_autor.png}
\end{center}
\end{frame}

\begin{frame}{Libro con Editor}
% Imagen: Ejemplo completo de referencia de libro con editor (Ed.) o (Eds.)
\begin{center}
\includegraphics[width=0.85\textwidth]{images/referencia_libro_editor.png}
\end{center}
\end{frame}

\begin{frame}{Libro Electrónico}
% Imagen: Ejemplo de libro electrónico con DOI o URL
\begin{center}
\includegraphics[width=0.85\textwidth]{images/referencia_libro_electronico.png}
\end{center}
\end{frame}

\begin{frame}{Capítulo de Libro Editado}
% Imagen: Ejemplo completo de capítulo de libro con páginas y editor
\begin{center}
\includegraphics[width=0.9\textwidth]{images/referencia_capitulo_libro.png}
\end{center}
\end{frame}

\begin{frame}{Libro Traducido}
% Imagen: Ejemplo de libro traducido con traductor y fechas original/traducción
\begin{center}
\includegraphics[width=0.85\textwidth]{images/referencia_libro_traducido.png}
\end{center}
\end{frame}

\section{Referencias de Artículos}

\begin{frame}{Anatomía de un Artículo Científico}
% Imagen: Primera página de artículo con elementos para referencia señalados
\begin{center}
\includegraphics[width=0.7\textwidth]{images/anatomia_articulo.png}
\end{center}
\end{frame}

\begin{frame}{Artículo de Revista Científica (Journal)}
% Imagen: Ejemplo completo de referencia de artículo con volumen, número, páginas
\begin{center}
\includegraphics[width=0.9\textwidth]{images/referencia_articulo_journal.png}
\end{center}
\end{frame}

\begin{frame}{Artículo con DOI}
% Imagen: Ejemplo de artículo con DOI en formato URL
\begin{center}
\includegraphics[width=0.85\textwidth]{images/referencia_articulo_doi.png}
\end{center}

\begin{block}{DOI}
Digital Object Identifier - identificador único y permanente
\end{block}
\end{frame}

\begin{frame}{Artículo en Línea sin DOI}
% Imagen: Ejemplo de artículo en línea con URL
\begin{center}
\includegraphics[width=0.85\textwidth]{images/referencia_articulo_url.png}
\end{center}
\end{frame}

\begin{frame}{Artículo de Periódico}
% Imagen: Ejemplo de artículo de periódico impreso y en línea
\begin{center}
\includegraphics[width=0.85\textwidth]{images/referencia_periodico.png}
\end{center}
\end{frame}

\begin{frame}{Artículo de Revista (Magazine)}
% Imagen: Ejemplo de artículo de revista de divulgación
\begin{center}
\includegraphics[width=0.85\textwidth]{images/referencia_magazine.png}
\end{center}
\end{frame}

\section{Tesis y Trabajos de Grado}

\begin{frame}{Tesis Publicada en Base de Datos}
% Imagen: Ejemplo completo de tesis en ProQuest u otra base de datos
\begin{center}
\includegraphics[width=0.9\textwidth]{images/referencia_tesis_bd.png}
\end{center}
\end{frame}

\begin{frame}{Tesis en Repositorio Institucional}
% Imagen: Ejemplo de tesis en repositorio con URL
\begin{center}
\includegraphics[width=0.85\textwidth]{images/referencia_tesis_repositorio.png}
\end{center}
\end{frame}

\section{Material Electrónico}

\begin{frame}{Página Web}
% Imagen: Ejemplo de referencia de página web con autor y sin autor
\begin{center}
\includegraphics[width=0.85\textwidth]{images/referencia_pagina_web.png}
\end{center}
\end{frame}

\begin{frame}{Entrada de Blog}
% Imagen: Ejemplo de entrada de blog
\begin{center}
\includegraphics[width=0.85\textwidth]{images/referencia_blog.png}
\end{center}
\end{frame}

\begin{frame}{Redes Sociales}
% Imagen: Ejemplos de Tweet, Facebook, Instagram
\begin{center}
\includegraphics[width=0.9\textwidth]{images/referencia_redes_sociales.png}
\end{center}
\end{frame}

\section{Material Audiovisual}

\begin{frame}{Película o Documental}
% Imagen: Ejemplo de referencia de película
\begin{center}
\includegraphics[width=0.85\textwidth]{images/referencia_pelicula.png}
\end{center}
\end{frame}

\begin{frame}{Serie de Televisión}
% Imagen: Ejemplo de referencia de serie de TV
\begin{center}
\includegraphics[width=0.85\textwidth]{images/referencia_serie_tv.png}
\end{center}
\end{frame}

\begin{frame}{Video de YouTube}
% Imagen: Ejemplo de referencia de video de YouTube
\begin{center}
\includegraphics[width=0.85\textwidth]{images/referencia_youtube.png}
\end{center}
\end{frame}

\begin{frame}{Podcast}
% Imagen: Ejemplo de referencia de podcast
\begin{center}
\includegraphics[width=0.85\textwidth]{images/referencia_podcast.png}
\end{center}
\end{frame}

\section{Referencias Legales}

\begin{frame}{Leyes y Decretos}
% Imagen: Ejemplo de referencia de ley colombiana
\begin{center}
\includegraphics[width=0.85\textwidth]{images/referencia_ley.png}
\end{center}
\end{frame}

\begin{frame}{Sentencias Judiciales}
% Imagen: Ejemplo de referencia de sentencia de Corte
\begin{center}
\includegraphics[width=0.85\textwidth]{images/referencia_sentencia.png}
\end{center}
\end{frame}

\begin{frame}{Tratados Internacionales}
% Imagen: Ejemplo de referencia de tratado internacional
\begin{center}
\includegraphics[width=0.85\textwidth]{images/referencia_tratado.png}
\end{center}
\end{frame}

\section{Conferencias e Informes}

\begin{frame}{Ponencia en Congreso}
% Imagen: Ejemplo de referencia de ponencia o conferencia
\begin{center}
\includegraphics[width=0.85\textwidth]{images/referencia_ponencia.png}
\end{center}
\end{frame}

\begin{frame}{Informe Gubernamental}
% Imagen: Ejemplo de referencia de informe de institución gubernamental
\begin{center}
\includegraphics[width=0.85\textwidth]{images/referencia_informe_gobierno.png}
\end{center}
\end{frame}

\section{Reglas Especiales}

\begin{frame}{Uso de "y" vs "\&"}
% Imagen: Ejemplos mostrando uso de "y" en español y "\&" en inglés
\begin{center}
\includegraphics[width=0.85\textwidth]{images/regla_y_ampersand.png}
\end{center}

\begin{block}{Adaptación al español}
En referencias en español: usar "y" sin coma previa
\end{block}
\end{frame}

\begin{frame}{Más de 20 Autores}
% Imagen: Ejemplo de referencia con más de 20 autores (primeros 19... último)
\begin{center}
\includegraphics[width=0.85\textwidth]{images/referencia_20_autores.png}
\end{center}
\end{frame}

\begin{frame}{Sin Autor}
% Imagen: Ejemplo de referencia iniciando con el título
\begin{center}
\includegraphics[width=0.85\textwidth]{images/referencia_sin_autor.png}
\end{center}
\end{frame}

\begin{frame}{Autor Corporativo}
% Imagen: Ejemplo de referencia con organización como autor
\begin{center}
\includegraphics[width=0.85\textwidth]{images/referencia_autor_corporativo.png}
\end{center}
\end{frame}

\begin{frame}{DOI vs URL}
% Imagen: Explicación visual de cuándo usar DOI y cuándo URL
\begin{center}
\includegraphics[width=0.85\textwidth]{images/doi_vs_url.png}
\end{center}
\end{frame}

\section{Formato de la Lista}

\begin{frame}{Orden Alfabético}
% Imagen: Ejemplo de lista de referencias en orden alfabético correcto
\begin{center}
\includegraphics[width=0.75\textwidth]{images/orden_alfabetico_referencias.png}
\end{center}
\end{frame}

\begin{frame}{Sangría Francesa}
% Imagen: Detalle de sangría francesa en referencias
\begin{center}
\includegraphics[width=0.8\textwidth]{images/sangria_francesa_detalle.png}
\end{center}
\end{frame}

\begin{frame}{Interlineado y Espaciado}
% Imagen: Ejemplo mostrando interlineado 2.0 sin espacio entre referencias
\begin{center}
\includegraphics[width=0.75\textwidth]{images/interlineado_referencias.png}
\end{center}
\end{frame}

\section{Errores Comunes}

\begin{frame}{Errores Frecuentes en Referencias}
% Imagen: Lista de los 10 errores más comunes con ejemplos
\begin{center}
\includegraphics[width=0.9\textwidth]{images/errores_comunes_referencias.png}
\end{center}
\end{frame}

\begin{frame}{Verificación: Citas vs Referencias}
\begin{alertblock}{Regla de Oro}
Toda cita debe tener su referencia \\
Toda referencia debe tener su cita
\end{alertblock}

% Imagen: Diagrama mostrando correspondencia entre citas y referencias
\begin{center}
\includegraphics[width=0.7\textwidth]{images/citas_referencias_correspondencia.png}
\end{center}
\end{frame}

\begin{frame}{Lista de Verificación: Referencias}
% Imagen: Checklist completo para revisar referencias
\begin{center}
\includegraphics[width=0.85\textwidth]{images/checklist_referencias.png}
\end{center}
\end{frame}

\section{Herramientas Útiles}

\begin{frame}{Gestores de Referencias}
\textbf{Software recomendado:}
\begin{itemize}
    \item Zotero (gratuito, código abierto)
    \item Mendeley (gratuito)
    \item EndNote (comercial)
    \item RefWorks (comercial)
\end{itemize}

% Imagen: Logos y capturas de pantalla de gestores de referencias
\begin{center}
\includegraphics[width=0.75\textwidth]{images/gestores_referencias.png}
\end{center}
\end{frame}

\begin{frame}{Recursos en Línea}
% Imagen: Lista de recursos útiles (APA Style, DOI.org, etc.) con URLs
\begin{center}
\includegraphics[width=0.85\textwidth]{images/recursos_online_apa.png}
\end{center}
\end{frame}

\begin{frame}{Sitio Oficial de APA Style}
% Imagen: Captura del sitio apastyle.apa.org
\begin{center}
\includegraphics[width=0.8\textwidth]{images/sitio_apa_style.png}
\end{center}

\begin{block}{Recurso Principal}
\url{https://apastyle.apa.org}
\end{block}
\end{frame}

\begin{frame}[standout]
\Huge ¿Preguntas?
\end{frame}

\begin{frame}{Resumen del Curso}
\begin{center}
\textbf{Hemos Visto:}
\vspace{0.5cm}

\begin{enumerate}
    \item Introducción a las Normas APA
    \item Estructura del trabajo académico
    \item Formato del documento
    \item Estilo de redacción
    \item Citas en el texto
    \item Lista de referencias
\end{enumerate}
\end{center}
\end{frame}

\begin{frame}{Práctica y Consulta}
\begin{block}{Recomendación Final}
\begin{itemize}
    \item Practique con trabajos reales
    \item Consulte el manual oficial ante dudas
    \item Use gestores de referencias
    \item Revise siempre antes de entregar
\end{itemize}
\end{block}

\vspace{0.5cm}
\begin{center}
\Large \textbf{¡Éxito en sus trabajos académicos!}
\end{center}
\end{frame}

\end{document}
